\subsection{The arithmetic degree of the original component lattice}
\label{sec:actual-component-section-degree}

We now apply the determinant theorem to the actual cusp lattices.
The geometric realization, its canonical line, the finite determinant
ideals and the Petersson norms are all retained.  In particular the
argument does not require integral lifts of arbitrary jets or the
selection of algebraic points of controlled height.

\paragraph{The original spaces.}
Fix one of the constructed components over its original totally complex
coefficient field $K$.  Let $X$ be its smooth projective toroidal full
flag compactification, $B_X$ its reduced boundary, and $U=X\setminus B_X$.
Write $D=\dim X\ge1$, and use the original automorphic lines
$\mathcal L(a)$ and canonical weight $\kappa_{\mathrm{can}}$.  Their
prescribed multiplicative identifications and
\[
                  \omega_X(B_X)=\mathcal L(\kappa_{\mathrm{can}})
\]
are those proved in Section~\ref{actual-geometric-positivity}.
Take the two original divided weights $u_1,u_2$ and the original integer
$T\ge D+1$.  Retain their original strict gap bounds
$\operatorname{gap}(u_i)\ge(D+3)^2+2r$, $i=1,2$, where
$r=p+q$ is the original rank.  Put
\[
 a_j=\kappa_{\mathrm{can}}+u_2+(D+j)u_1,\qquad 0\le j\le T.
\]
For any weight $a$ used here, let $C(a)$ be its original $K$-space of
cuspidal sections supported on this component, with its original
integral lattice $M(a)$ and component Petersson norms
$\|\cdot\|_{\sigma,a}$, for all embeddings $\sigma:K\hookrightarrow
\mathbb C$.  These norms use the original component probability
measure.

Choose the full integral frame of $C(u_1)$ constructed in
Section~\ref{shared-actual-integral-cusp-frame}.  Denote its generators
by $w_1,\ldots,w_n$, its dimension by $N$, its generator bound by
$\Lambda\ge1$, and its upper Bergman bound by $\beta\ge1$.
Thus the $\mathcal O_K$-span of the $w_i$ is exactly $M(u_1)$,
$1\le N\le n$, and
\begin{equation}\label{eq:actual-section-frame}
 \|w_i\|_{\sigma,u_1}\le\Lambda,\qquad
 \|s(x)\|^2\le\beta\|s\|_{\sigma,u_1}^2.
\end{equation}
Let $\Lambda_0,\Lambda_D,\Lambda_T\ge1$ be the supplied short
spanning bounds for $M(a_0),M(a_D),M(a_T)$.

Write $\delta$ and $h$ for the original image degree and normalized
image height of the system $C(u_1)$ on this component.  These use
the full lattice $M(u_1)$ and its stated Hermitian norms.  The precise
height formula is recalled and verified below.  Put
\begin{gather*}
 d=\dim_K C(a_T),\qquad
 q=\binom{T-1}{D}\delta,\qquad \vartheta=\frac qd,\\
 c=\log(n\Lambda)+\frac{\log n+1}{2},\qquad
 a_n=\frac{\log n+1}{2}.
\end{gather*}

\begin{theorem}[The original component section estimate]
\label{thm:actual-component-section-degree}
For the constructed geometric data and the original weight-size
conditions, one has $d\ge q>0$ and
\begin{equation}\label{eq:actual-component-section-strong}
\begin{split}
\frac{\widehat{\deg} M(a_T)}d\ge{}&
 \vartheta\left(Th-Ta_n-D(D+1)c-\log\Lambda_0
                                      -\frac T2\log\beta\right)\\
&-(1-\vartheta)\left(Tc+\frac T2\log\beta
                             +\log\Lambda_D+\log\Lambda_T\right).
\end{split}
\end{equation}
Arithmetic degree is divided by $[K:\mathbb Q]$ and is the original
component degree.

In particular, set
\[
 \rho=\log\Lambda+
 26\log\!\left(nD\delta(2+|h|+\log\Lambda)\right),
 \qquad G=\frac{(D+1)(\log n+1)}2+1,\qquad z=1-\vartheta .
\]
Then $0\le z\le1$ and the literal numerical estimate required by
the component construction holds:
\begin{equation}\label{eq:actual-component-section-original}
\begin{split}
\frac{\widehat{\deg}M(a_T)}d\ge{}&
(1-z)\bigl(Th-3(D+1)|h|-\log\Lambda_0
 -26D^2\log(e^\rho T)\\[-2pt]
&\hspace{36mm}-(T+D+1)\log\beta-(T+D^2)(G+2)\bigr)\\
&-z\bigl(T\log(2e^\rho\sqrt\beta)
                              +\log\Lambda_D+\log\Lambda_T\bigr).
\end{split}
\end{equation}
At the original family size $P\ge0$ and weight size $\kappa\ge0$, the already
constructed short families give
\[
 \Lambda_0=\exp\!\left(DP(\kappa/T+1)\right),\qquad
 \Lambda_D=\Lambda_T=\exp(P\kappa).
\]
Thus the original remainder coefficient is exactly
$J=DP(\kappa/T+1)$.  The theorem supplies the component arithmetic
section inequality; the separate all-point height comparison and
any still unconstructed numerical family are not assumptions
silently discharged by this application.
\end{theorem}

\begin{proof}
\emph{The cusp space and the canonical extension.}
For each weight under consideration there is an equality over $K$
\begin{equation}\label{eq:actual-section-cusp-extension}
             C(a)=H^0(X,\mathcal L(a)(-B_X)).
\end{equation}
Here the line on the right is the subcanonical extension.  The
analytic--algebraic toroidal comparison identifies its automorphic
bundle and Fourier--Jacobi coordinates with those of the original
moduli problem; the subcanonical comparison identifies its sections
with holomorphic cuspidal forms
\cite{positive-lan-analytic,positive-lan-stroh}.
The passage from the automorphic bundle to its full flag line is the
usual highest-weight realization, with the same ordered characters.
The original weights satisfy the needed dominance conditions.

We spell out why no growth condition changes this space.  In a smooth
toroidal boundary chart, a canonical Hodge frame and its inverse have
at most a fixed power of the logarithmic cusp coordinates as growth.
The same holds for every fixed tensor and flag character.  Rapid decay
of a holomorphic cuspidal form therefore makes its coefficient tend
to zero at a generic point of each boundary divisor
\cite{positive-cusp-decay}.  Its algebraic meromorphic coefficient
has positive order there.  Smoothness, hence normality, extends it
across the remaining codimension at least two subset.  It is a
section of the negative reduced boundary.  Conversely such a section
has vanishing constant Fourier--Jacobi coefficients, so its
automorphic reading is cuspidal; the same rapid-decay theorem gives
the original bounded Petersson reading.  The forms have the
prescribed holomorphic highest weight and hence finite infinitesimal
character, as in the earlier cusp verification.  Finally the
identifications are algebraic over the given $K$, and cohomology
base change descends their equality from an embedding into
$\mathbb C$.  This is an equality of rational spaces, and does not
replace their integral lattices.

\emph{The exact geometric realization.}
Principalize the base ideals of $C(u_1)$ and $C(u_2)$
simultaneously over $K$, using their prescribed birational
modification $\pi:Y\to X$~\cite{WlodarczykResolution}.
One can principalize their product; the two factors of an invertible
product ideal are invertible fractional ideals, so both original
ideals become invertible.  Both systems are generated on $U$, so
the modification is an isomorphism there.  The constructed $X$
is smooth projective, and the principalization has smooth projective
source $Y$.  Let $F_1,F_2$ be the actual effective fixed divisors of
the pulled-back systems.  They are supported outside $U$.  Set
\begin{equation}\label{eq:actual-section-lines}
 L=\pi^*\mathcal L(u_1)(-F_1),\qquad
 L_2=\pi^*\mathcal L(u_2)(-F_2),\qquad
 A=\pi^*\mathcal L(u_2+Du_1).
\end{equation}
The two divided systems generate $L,L_2$ everywhere and are
birational.  Thus $L,L_2$ are nef and big.  The original weight gaps
make $\mathcal L(u_2+Du_1)$ nef and big by
Theorem~\ref{actual-positive-cusp-theorem}; pullback under $\pi$
preserves both properties.  This proves the required assertion for
$A$.  The canonical section with divisor $F_2+DF_1$ is a section
\[
            e\in H^0(Y,A\otimes L_2^{-1}\otimes L^{-D})
\]
which is nonzero throughout $U$.

Put $\mathcal G_j=\omega_Y\otimes A\otimes L^j$.
Use the prescribed Kodaira--Spencer canonical identification on $U$
and the actual sections of the $F_i$ to identify $C(a_j)$ with
a space $V_j\subset H^0(U,\mathcal G_j)$.
For clarity the needed inclusions follow without a vanishing
theorem for higher direct images.  One has
\begin{equation}\label{eq:actual-section-canonical-pushforward}
                         \pi_*\omega_Y=\omega_X.
\end{equation}
Indeed a proper birational morphism to the smooth $X$ is an
isomorphism above the generic point of every divisor.  A regular
top form on $Y$ is therefore a regular top form on $X$ outside
codimension at least two, and extends uniquely because $\omega_X$
is locally free and $X$ is normal.  Conversely the pullback of a
regular top form is regular: locally its coefficients are multiplied
by the Jacobian determinant of a morphism between smooth varieties.
These two operations are inverse on a dense open.

By the projection formula, \eqref{eq:actual-section-canonical-pushforward}
and \eqref{eq:actual-section-cusp-extension},
\[
 H^0(Y,\omega_Y A)
   =H^0(X,\omega_X\mathcal L(u_2+Du_1))=C(a_0).
\]
Thus $V_0=H^0(Y,\mathcal G_0)$.  Multiplication by the section
of $jF_1$ embeds
$\mathcal G_j$ into
$\omega_Y\otimes\pi^*\mathcal L(u_2+(D+j)u_1)$.
Its global sections descend, by the same argument, to $C(a_j)$.
Hence
\[
                    H^0(Y,\mathcal G_j)\subset V_j
                         \quad(0\le j\le T).
\]
Multiplication by the divided section of $L$ corresponds exactly
to multiplication by the original cusp section of weight $u_1$:
the additional factor of $F_1$ is cancelled by the next division.
These are precisely the spaces and products in
Proposition~\ref{prop:geometric-section-filtration}.  Its finite
section and Koszul argument supplies all of the required geometric
filtration, including the surjection at degree zero and the equality
at the deepest step.  No arithmetic estimate on $e$ or the $F_i$
enters that argument.

\emph{The original arithmetic data.}
Transport $M(a_j)$ and its Hermitian norms through the preceding
$K$-linear identification.  This is exactly how the original
component section system assigns its lattice and metric
\cite{actual-section-source}.  Thus lattice membership and every
archimedean norm are equal before and after transport.
In particular no lattice of all integral global sections on $Y$
is substituted for $M(a_j)$.

Multiplication of two original integral cusp sections remains
integral at every original integral test object, and preserves the
component support.  The exact weight identity
$u_1+a_j=a_{j+1}$ proves the required lattice stability.
Pointwise metric multiplicativity and the original probability
measure, together with \eqref{eq:actual-section-frame}, give
\[
 \|s\,v\|_{\sigma,a_{j+1}}^2
   \le \|s\|_{\sup,\sigma,u_1}^2\|v\|_{\sigma,a_j}^2
   \le\beta\|s\|_{\sigma,u_1}^2\|v\|_{\sigma,a_j}^2.
\]
The short families at $0,D,T$ transport without a change in bound.

There is also no factor of two hidden in the degree.  If
$a=[K:\mathbb Q]$ and the covolume uses one complex embedding
per conjugate pair with its underlying real Euclidean measure,
then for a rank-$d$ full Hermitian lattice $M$
\[
 \widehat{\deg}M
   =-\frac1a\log\operatorname{covol}(M)
      -\frac d2\log2+\frac d2\log\operatorname{rd}(K).
\]
To verify this, express $\bigwedge^dM$ by a fractional ideal
$I$ in any $K$-basis.  Its real covolume is
$2^{-ad/2}|\Delta_K|^{d/2}\operatorname NI$
times the product of the Hermitian Gram determinants over the
complex places.  Substitution gives
$-a^{-1}(\log\operatorname NI+
\tfrac12\sum_\sigma\log\det\operatorname{Gram}_\sigma)$,
the normalized arithmetic degree used in the determinant theorem.
The original squared-norm slope is twice this degree divided by
rank.  Hence the degree just transported is the stated original
component degree.

\emph{The exact image height.}
Use any $K$-basis of $C(u_1)$, and let $f$ be the irreducible Chow
form of its image, of degree $\delta$ in each of $D+1$ blocks.
Its finite height is the Gauss height after substituting the full
integral generators $w_i$.  Its archimedean height is the mean of
$\log|f|$ over the product of the intrinsic Hermitian unit spheres.
Denote their normalized sum by $H(f)$, exactly as in
Section~\ref{sec:arithmetic-determinant-degree}.  The original image
height satisfies
\begin{equation}\label{eq:actual-section-height}
       h=\frac{H(f)}{(D+1)\delta}
                   +\frac12\sum_{i=1}^{N-1}\frac1i .
\end{equation}
This is its defining Chow-height normalization, not an estimate on
selected algebraic points.

To check independence of the actual frame, two finite generating
families of the same full lattice express one another by integral
matrices.  Integral linear substitution cannot increase a finite
Gauss norm, so the two opposite substitutions prove equality at
every finite place.  At infinity the change between the two coordinate
realizations, each with its transported Hermitian Gram matrix, is
an isometry of the same Hermitian space.  It preserves the spherical
probability measure; the basis-change matrix need not be unitary
for an unrelated standard form.  Chow forms in two bases correspond by that
linear substitution, up to a nonzero scalar; the product formula
cancels the scalar.  This proves \eqref{eq:actual-section-height}
for the original frame and the original $h$, without a projective
maximum norm or an extra period factor.  Birationality identifies
the image degree with the degree of the good finite section in
the geometric filtration.  Since $N\le n$,
\[
        \frac{H(f)}{(D+1)\delta}\ge h-a_n.
\]
All hypotheses of Theorem~\ref{thm:arithmetic-determinant-degree}
are now supplied for the original component spaces and lattices.
It gives \eqref{eq:actual-component-section-strong} and $d\ge q>0$.

\emph{The literal numerical remainder.}
The argument of the logarithm in $\rho-\log\Lambda$ is at least
$2n$.  As $n,\Lambda\ge1$, this implies $0\le c\le\rho$.
Also $D(D+1)\le2D^2$ and $T\ge2$, so
\[
 D(D+1)c\le26D^2(\rho+\log T),\qquad
 \frac T2\log\beta\le(T+D+1)\log\beta,\qquad
 Ta_n\le(T+D^2)(G+2).
\]
The added $3(D+1)|h|$ is nonnegative.  The complementary term
is bounded because
$c+\tfrac12\log\beta\le
\log(2e^\rho\sqrt\beta)$.
Multiplication by the nonnegative factors $\vartheta$ and
$1-\vartheta$ now proves \eqref{eq:actual-component-section-original}.
Finally Corollary~\ref{cor:actual-hs-short-families} supplies the displayed
$\Lambda_0,\Lambda_D,\Lambda_T$ in the original family, while
the full frame theorem supplies $n,\Lambda,\beta$ with their
original bounds.  This proves the stated value of $J$.
\end{proof}

The stronger estimate also combines linearly over components.
If $d_Z,q_Z,h_Z$ are the quantities above, put
$d_{\mathrm{tot}}=\sum_Zd_Z$,
$\Theta=\sum_Zq_Z/d_{\mathrm{tot}}$ and
$\bar h=(\sum_Zq_Zh_Z)/(\sum_Zq_Z)$.
With common upper bounds for the logarithmic costs, summation gives
\[
\begin{split}
\frac{\sum_Z\widehat{\deg}M_Z(a_T)}{d_{\mathrm{tot}}}
\ge{}&\Theta T\bar h-\frac T2\log\beta\\
&-\Theta\bigl(Ta_n+D(D+1)c+\log\Lambda_0\bigr)\\
&-(1-\Theta)\bigl(Tc+\log\Lambda_D+\log\Lambda_T\bigr).
\end{split}
\]
All dependence on heights is through the original signed weighted
mean.  Bounds for that mean, the rank ratios and the remaining
original numerical family are still required at the final
geometric deduction.
